% !TEX root = ../main.tex
%-----------------------------------------------------------------------
%  Title block, the status warning, and the question.
%  \input by ../main.tex -- not compiled on its own.
%-----------------------------------------------------------------------
\begin{center}
{\large\bfseries (Tanvir Testing from Sverdrup)Scoping Note: DINO at 1/4\dg}\\[2pt]
{\small Moving the adjoint programme from a parameterised-eddy basin
 to a partly eddy-resolving one}\\[4pt]
{\small Tanvir Shahriar \quad$\cdot$\quad DINO configuration, MITgcm
 \quad$\cdot$\quad \today}
\end{center}

\hd{Status: placeholder.}
Nothing is written, nothing is configured, and no setup exists. This note is
here so the direction has somewhere to accumulate, not because any of it has
been thought through. Every figure below is arithmetic applied to measurements
taken at 1\dg, never a measurement of a 1/4\dg\ configuration --- none has been
run. Read the numbers as an order of magnitude that decides whether the
direction is worth planning properly, and nothing more.

\hd{The question.}
Everything in this repository is the 1\dg\ basin: the 200-year spin-up, the
reference adjoint, and the perturbation ensemble the surrogate direction is
built on. At 1/4\dg\ the flow is eddy-permitting, so the state the adjoint
linearises about is a different one --- more energetic, and carrying variability
the coarse configuration cannot produce. That raises three things, none of them
answered here.

\smallskip
\begin{tabularx}{\textwidth}{@{}Y{0.42}Y{1.58}@{}}
\toprule
\textbf{Question} & \textbf{Why it matters} \\
\midrule
Does the sensitivity structure survive the resolution change?
  & If the 1\dg\ sensitivity maps are a good guide to the 1/4\dg\ ones, the
    coarse configuration is a cheap laboratory for the fine one. If they are
    not, that is its own result. \\
\addlinespace
Does the adjoint stay usable over a five-year sweep?
  & A more energetic, less dissipative flow is the regime where long adjoints
    are known to misbehave. \code{adjViscBoost} may go from a variant to a
    requirement. \\
\addlinespace
What does it cost?
  & The constraint that shapes everything else, and the only one of the three
    that can be estimated before anything is built. Section~2. \\
\bottomrule
\end{tabularx}

\hd{What carries over unchanged.}
Most of the machinery is resolution-agnostic: the Tapenade build path and the
\code{genmake2} patch, \code{tools/machine\_env.sh}, \code{tools/submit.sh}, the
\code{code\_tap/}~+~\code{input\_tap/} split, the namelist-variant convention,
and the DINO analysis notebooks, which take their shape from the run directory
through \code{xmitgcm.open\_mdsdataset(read\_grid=True)} rather than hardcoding
one. Two things would need editing by hand: the SOMA notebooks carry
\code{sNx}/\code{sNy}/\code{OLx}/\code{OLy}/\code{Nr} as literals, and the
viscosity-binary notebook hardcodes \code{Nx, Ny, Nr = 51, 198, 36}.
